% GENERATED FILE. Edit canonical lessons, then run npm run generate:books.
% canonical-source-hash: fnv1a64:bae9e01178168085
% canonical-lessons: JA-C90-odoroku, JA-C90-todoku, JA-C90-itadaku, JA-C90-kayou, JA-C90-tsutsumu

\chapter{To be surprised, To arrive, To receive, To go regularly, To wrap}
\label{ch:ja-to-be-surprised-to-arrive-to-receive-to-go-regularly-to-wrap}
\hlchaptermodality{90}

\begin{chapteropening}
\textbf{By the end of this chapter:} \emph{I can say to be surprised, to arrive, to receive, to go regularly and to wrap.}

Complete the last lesson of chapter 90: i can say to be surprised, to arrive, to receive, to go regularly and to wrap.
\end{chapteropening}

\section[\texorpdfstring{\ja{おどろく} (odoroku)}{おどろく (odoroku)}]{\ja{おどろく} (odoroku) — to be surprised}

\label{lesson:JA-C90-odoroku}



\begin{quote}
Before the new one: say the Japanese for to grab, then the Japanese for to fold.
\end{quote}

\subsection*{You'll want to know: \ja{おどろく}}
\textbf{\ja{おどろく}} — \emph{odoroku} — \textquotedblleft{}to be surprised\textquotedblright{}.

\begin{grammarlens}[title={What you've built}]
One more word for this chapter.
\end{grammarlens}

\subsection*{Your turn}
\begin{itemize}
\raggedright
  \item \emph{Say it:} \emph{odoroku}
  \item \emph{Say it:} \emph{odoroku}, once more
  \item \emph{Say it:} say \emph{odoroku}
  \item \emph{Recall:} say the Japanese for to grab, then the Japanese for to fold, then say \emph{odoroku} again
\end{itemize}

\begin{tcolorbox}[breakable,colback=teal!4,colframe=teal!35!black,title={Before you move on}]
What does \ja{おどろく} mean? (\textquotedblleft{}To be surprised\textquotedblright{}.) Say it once more.
\end{tcolorbox}
\section[\texorpdfstring{\ja{とどく} (todoku)}{とどく (todoku)}]{\ja{とどく} (todoku) — to arrive, to reach}

\label{lesson:JA-C90-todoku}



\begin{quote}
Before the new one: say the Japanese for to fold, then the Japanese for to be surprised.
\end{quote}

\subsection*{You'll want to know: \ja{とどく}}
\textbf{\ja{とどく}} — \emph{todoku} — \textquotedblleft{}to arrive, to reach\textquotedblright{}.

\begin{grammarlens}[title={What you've built}]
One more word for this chapter.
\end{grammarlens}

\subsection*{Your turn}
\begin{itemize}
\raggedright
  \item \emph{Say it:} \emph{todoku}
  \item \emph{Say it:} \emph{todoku}, once more
  \item \emph{Say it:} say \emph{todoku}
  \item \emph{Recall:} say the Japanese for to fold, then the Japanese for to be surprised, then say \emph{todoku} again
\end{itemize}

\begin{tcolorbox}[breakable,colback=teal!4,colframe=teal!35!black,title={Before you move on}]
What does \ja{とどく} mean? (\textquotedblleft{}To arrive, to reach\textquotedblright{}.) Say it once more.
\end{tcolorbox}
\section[\texorpdfstring{\ja{いただく} (itadaku)}{いただく (itadaku)}]{\ja{いただく} (itadaku) — to receive (politely)}

\label{lesson:JA-C90-itadaku}



\begin{quote}
Before the new one: say the Japanese for to be surprised, then the Japanese for to arrive.
\end{quote}

\subsection*{You'll want to know: \ja{いただく}}
\textbf{\ja{いただく}} — \emph{itadaku} — \textquotedblleft{}to receive (politely)\textquotedblright{}.

\begin{grammarlens}[title={What you've built}]
One more word for this chapter.
\end{grammarlens}

\subsection*{Your turn}
\begin{itemize}
\raggedright
  \item \emph{Say it:} \emph{itadaku}
  \item \emph{Say it:} \emph{itadaku}, once more
  \item \emph{Say it:} say \emph{itadaku}
  \item \emph{Recall:} say the Japanese for to be surprised, then the Japanese for to arrive, then say \emph{itadaku} again
\end{itemize}

\begin{tcolorbox}[breakable,colback=teal!4,colframe=teal!35!black,title={Before you move on}]
What does \ja{いただく} mean? (\textquotedblleft{}To receive (politely)\textquotedblright{}.) Say it once more.
\end{tcolorbox}
\section[\texorpdfstring{\ja{かよう} (kayou)}{かよう (kayou)}]{\ja{かよう} (kayou) — to go regularly}

\label{lesson:JA-C90-kayou}



\begin{quote}
Before the new one: say the Japanese for to arrive, then the Japanese for to receive.
\end{quote}

\subsection*{You'll want to know: \ja{かよう}}
\textbf{\ja{かよう}} — \emph{kayou} — \textquotedblleft{}to go regularly\textquotedblright{}.

\begin{grammarlens}[title={What you've built}]
One more word for this chapter.
\end{grammarlens}

\subsection*{Your turn}
\begin{itemize}
\raggedright
  \item \emph{Say it:} \emph{kayou}
  \item \emph{Say it:} \emph{kayou}, once more
  \item \emph{Say it:} say \emph{kayou}
  \item \emph{Recall:} say the Japanese for to arrive, then the Japanese for to receive, then say \emph{kayou} again
\end{itemize}

\begin{tcolorbox}[breakable,colback=teal!4,colframe=teal!35!black,title={Before you move on}]
What does \ja{かよう} mean? (\textquotedblleft{}To go regularly\textquotedblright{}.) Say it once more.
\end{tcolorbox}
\section[\texorpdfstring{\ja{つつむ} (tsutsumu)}{つつむ (tsutsumu)}]{\ja{つつむ} (tsutsumu) — to wrap}

\label{lesson:JA-C90-tsutsumu}



\begin{quote}
Before the new one: say the Japanese for to receive, then the Japanese for to go regularly.
\end{quote}

\subsection*{You'll want to know: \ja{つつむ}}
\textbf{\ja{つつむ}} — \emph{tsutsumu} — \textquotedblleft{}to wrap\textquotedblright{}.

\begin{grammarlens}[title={What you've built}]
One more word for this chapter.
\end{grammarlens}

\subsection*{Your turn}
\begin{itemize}
\raggedright
  \item \emph{Say it:} \emph{tsutsumu}
  \item \emph{Say it:} \emph{tsutsumu}, once more
  \item \emph{Say it:} say \emph{tsutsumu}
  \item \emph{Recall:} say the Japanese for to receive, then the Japanese for to go regularly, then say \emph{tsutsumu} again
\end{itemize}

\begin{tcolorbox}[breakable,colback=teal!4,colframe=teal!35!black,title={Before you move on}]
What does \ja{つつむ} mean? (\textquotedblleft{}To wrap\textquotedblright{}.) Say it once more.
\end{tcolorbox}
